\chapter{Detección de comportamiento con modelos ocultos de Markov}
\label{ch:o3}

El capítulo anterior cerró señalando dos límites de la cadena de Markov visible que
su formulación no podía salvar: no distinguía el régimen de comportamiento de cada
día (si el ave permanecía en su zona o estaba en plena migración) ni la ruta propia
de cada individuo. Este capítulo aborda el primero. La cuestión que plantea es si un
modelo no supervisado puede descubrir, en las series de posiciones diarias de
\textit{Larus fuscus}, esos regímenes que la cadena visible promediaba sin
separarlos: el movimiento de la especie alterna fases sedentarias (en las zonas de
cría e invernada) con fases de desplazamiento dirigido (los pasos migratorios de
primavera y otoño), pero el conjunto de datos no trae ninguna etiqueta que indique,
para un día concreto, en cuál de los dos regímenes se encontraba el ave. El resultado
no será todavía una predicción de la posición, sino una etiqueta de comportamiento
por ave y día.

La cadena de Markov visible no podía separar esos regímenes porque su único estado,
la celda ocupada, no distinguía si el ave permanecía en su zona o migraba. El modelo
oculto de Markov sí los separa al introducir una variable latente, el estado de
comportamiento, del que depende el movimiento de cada día. Su marco formal se expone
en el capítulo \ref{ch:preliminares} (sección \ref{sec:prelim-hmm}); aquí se ajusta un
modelo con dos estados ocultos, que se nombran \emph{estacionario} (movimiento local) y
\emph{migración} (desplazamiento dirigido de largo alcance), con la justificación de su
número en la sección \ref{sec:o3-modelo}.

El modelo es global: sus parámetros los comparten las 82 aves del conjunto. A
diferencia del capítulo anterior, donde las rutas individuales hacían que un modelo
poblacional generalizara mal a un ave no vista, aquí el estado se infiere a partir de
variables cinemáticas (la distancia recorrida y el cambio de rumbo, sección
\ref{sec:o3-features}) comparables entre individuos, sobre las que sí es razonable
compartir parámetros.

Partiendo de la misma tabla diaria del capítulo \ref{ch:o1} (una posición por ave y
día natural, con los huecos representados de forma explícita), el capítulo produce una
salida de doble uso: cada observación (ave, día) recibe una etiqueta de estado y la
probabilidad de cada régimen, que los modelos supervisados del capítulo \ref{ch:o4}
consumen como variable de entrada de alto nivel y que, a la vez, permite estratificar
por régimen el análisis del error del capítulo \ref{ch:o5}, separando el acierto de
los días sedentarios del de los días de migración. El estado detectado es la pieza que
conecta la fenología de la especie con los modelos predictivos del resto del trabajo.

\section{Variables del movimiento}
\label{sec:o3-features}

El modelo no observa la posición directamente, sino un vector de variables que resume
el movimiento de cada día y que es la observación $O_t$ que emite el estado oculto
(capítulo \ref{ch:preliminares}, ecuación \ref{eq:prelim-emision}). La elección de esas
variables determina qué puede separar el modelo: deben distinguir un día sedentario de
uno de migración a partir de la geometría de la trayectoria, no de la posición
absoluta. Se emplean dos variables cinemáticas, calculadas sobre las posiciones diarias
del capítulo \ref{ch:o1}.

La primera es el desplazamiento diario, \texttt{step\_in\_km}: la distancia de
Haversine (ecuación \ref{eq:haversine}) entre la posición de ayer y la de hoy, en
kilómetros y sin transformación. Es la señal primaria del régimen, porque separa por
sí sola los dos modos del movimiento de la especie: un día sedentario recorre unos
pocos kilómetros y uno de migración, cientos. Conservar la escala original en lugar de
comprimirla con un logaritmo es una decisión deliberada que la sección
\ref{sec:o3-modelo} justifica. La segunda es el cambio de rumbo,
\texttt{cos\_turning\_in}: el coseno del ángulo entre el rumbo de hoy (el del tramo de
ayer a hoy) y el del día anterior. Distingue la forma del movimiento: un valor próximo
a $+1$ indica rumbo sostenido (trayectoria rectilínea, propia del vuelo migratorio
dirigido), uno próximo a $-1$ una inversión de la marcha y uno próximo a $0$ un giro de
90\textdegree. Se codifica como coseno, y no como el ángulo en valor absoluto, porque
es una función suave en toda la recta que se ajusta mejor a la emisión gaussiana del
modelo, mientras que el valor absoluto introduce en el cero un quiebro que distorsiona
la distribución.

Ambas variables son entrantes: describen el tramo ya recorrido hasta el día $t$ (de
$t-1$ a $t$), no el que vendrá después (de $t$ a $t+1$). Así caracterizan el día con
información disponible en el momento y no con la posición futura, una propiedad que la
sección \ref{sec:o3-modelo} retoma al inferir el estado.

Estas dos variables forman el vector de observación de un primer modelo, el cinemático
(modelo A). Un segundo modelo (modelo B) añade tres variables de contexto ambiental y
temporal: la cobertura de vegetación baja (\texttt{veg\_low}) y alta
(\texttt{veg\_high}), tomadas de las covariables ambientales del conjunto Movebank, y
las horas de luz (\texttt{daylight\_hours}), calculadas con una fórmula astronómica a
partir de la latitud y la fecha. Estas tres variables se incorporan para tratar de
reflejar mejor la realidad biológica de los patrones migratorios, que no dependen solo
del movimiento sino también del hábitat y de la época del año. Las dos coberturas de
vegetación se mantienen separadas, porque la baja (pastos, cultivos) y la alta (bosque)
son biomas distintos que sumar en una sola variable confundiría. La comparación entre ambos modelos se resuelve en la sección
\ref{sec:o3-resultados}; hasta entonces, el que incorpora el contexto es el de
referencia.
